% ====================================================================
% Capítulo 3: Mecanismos estructurales en agua líquida
% ====================================================================
\chapter{Mecanismos estructurales en agua líquida: enlaces de hidrógeno y supersolidez}
\label{ch:mecanismos_agua}

\section{El paradigma de Zhang \emph{et al.}: memoria de enlaces de hidrógeno y supersolidez de piel}

Zhang \emph{et al.}~\cite{zhang2014} propusieron una explicación cuantitativa del 
efecto Mpemba basada en dos fenómenos íntimamente acoplados:

\begin{enumerate}
\item La \textbf{memoria de los enlaces de hidrógeno} ($\mathrm{O:H{-}O}$): el 
intervalo entre el átomo de oxígeno donador y el aceptor posee una dinámica 
relajacional cuyo tiempo característico depende fuertemente de la temperatura 
inicial $\Tinit$ de la cual fue templada la muestra.
\item La \textbf{supersolidez de piel del agua}: la región superficial 
(\SIrange{1}{2}{\milli\meter} de espesor) posee densidad reducida y conductividad 
térmica aumentada respecto al volumen, modificando localmente el campo de 
temperatura.
\end{enumerate}

\subsection{Estructura del enlace $\mathrm{O:H{-}O}$ y sus dos componentes}

A diferencia de modelos clásicos que representan el enlace de hidrógeno como un 
contacto único, Zhang \emph{et al.} lo descomponen en dos sub-enlaces:

\begin{itemize}
\item Un \textbf{enlace covalente $\mathrm{H{-}O}$} de \SI{0.1}{\nano\meter} y 
energía $E_{H} \sim \SI{4.0}{\electronvolt}$.
\item Un \textbf{no-enlace electrostático $\mathrm{O{:}H}$} de 
\SI{0.18}{\nano\meter} y energía $E_{L} \sim \SI{0.1}{\electronvolt}$.
\end{itemize}

\noindent Ambos están \emph{acoplados} por la repulsión Coulombiana entre los 
electrones del oxígeno (\textit{Coulomb repulsion}, C-repulsion): cualquier cambio 
en $d_{\mathrm{H{-}O}}$ induce un cambio en $d_{\mathrm{O{:}H}}$ en sentido 
opuesto. Esto da al sistema una \emph{histéresis} interna que actúa como memoria 
térmica.

Bajo calentamiento, $d_{\mathrm{H{-}O}}$ se acorta (anómalamente) y 
$d_{\mathrm{O{:}H}}$ se alarga; durante el enfriamiento ocurre lo opuesto, pero 
con una constante temporal $\tau_{\mathrm{HB}}$ que depende de la historia térmica.

\subsection{Ecuación de Fourier modificada}

El modelo de Zhang \emph{et al.} discretiza el agua en una columna 
unidimensional de longitud $L$ y resuelve la ecuación de Fourier:

\begin{equation}
    \rho(\theta) \, c_{p}(\theta) \, \pdv{\theta}{t} 
    = \pdv{x} \left[ \kappa(\theta,x) \pdv{\theta}{x} \right] 
    - \rho c_{p} v_{x}(x) \, \pdv{\theta}{x} 
    + q_{\mathrm{mem}}(x,t),
    \label{eq:fourier_zhang}
\end{equation}
%
donde $\theta$ es la temperatura (en \si{\celsius}), $\rho$, $c_{p}$ y $\kappa$ 
son funciones del campo local de temperatura, $v_{x}$ es la velocidad 
convectiva y $q_{\mathrm{mem}}$ representa la fuente/sumidero térmico asociado a 
la relajación de la memoria de los enlaces de hidrógeno:
\begin{equation}
    q_{\mathrm{mem}}(x,t) = C_{\mathrm{mem}} \, \dv{d_{\mathrm{OH}}}{t}, 
    \qquad
    \dv{d_{\mathrm{OH}}}{t} = - \frac{d_{\mathrm{OH}} - d_{\mathrm{OH}}^{\mathrm{eq}}(\theta)}{\tau_{\mathrm{HB}}(\Tinit)},
    \label{eq:hb_memory}
\end{equation}
%
con tiempo característico
\begin{equation}
    \tau_{\mathrm{HB}}(\Tinit) = \tau_{0} \, \exp\!\left[-\frac{\Tinit - \theta_{0}}{T_{\mathrm{scale}}}\right].
    \label{eq:tau_hb}
\end{equation}
%
La \emph{decreciente} dependencia $\tau_{\mathrm{HB}}(\Tinit) \downarrow$ con 
$\Tinit \uparrow$ implica que el agua más caliente \emph{libera su memoria más 
rápido} al ser puesta en contacto con el baño frío. Es esta asimetría la que 
genera, en el modelo de Zhang, la disipación más eficiente de la entalpía total 
desde la muestra inicialmente caliente.

\subsection{Supersolidez de piel}

La hipótesis adicional ---ya documentada por Zhang en trabajos previos--- es que 
la región superficial del agua, donde las moléculas tienen coordinación 
incompleta, exhibe propiedades de \emph{supersólido}: 

\begin{itemize}
\item Densidad reducida: $\rho_{\mathrm{skin}}/\rho_{\mathrm{bulk}} \approx 0.75$.
\item Conductividad térmica aumentada: $\kappa_{\mathrm{skin}}/\kappa_{\mathrm{bulk}} \approx 4/3$.
\item Frecuencia fonónica O--H alta ($\sim \SI{3450}{\per\centi\meter}$, frente a 
$\SI{3200}{\per\centi\meter}$ en el bulk).
\end{itemize}

\noindent Esta capa actúa como un \emph{drenaje térmico} preferente para la 
disipación de calor durante el enfriamiento. La \cref{eq:fourier_zhang} debe 
resolverse con condiciones de frontera de tipo Robin en la interfase 
piel-ambiente:
\begin{equation}
    h_{\mathrm{drain}} \left[\theta(0,t) - \theta_{\mathrm{bath}}\right] 
    = \kappa_{\mathrm{skin}} \pdv{\theta}{x}\bigg|_{x=0}.
    \label{eq:robin_bc}
\end{equation}

\section{Mecanismos microscópicos: el modelo de Jin y Goddard (2015)}

Mientras Zhang \emph{et al.} trabajan a nivel macroscópico (campo $\theta(x,t)$), 
Jin y Goddard~\cite{jin_goddard2015} investigaron el origen molecular del efecto 
mediante simulaciones de dinámica molecular clásica.

\subsection{Protocolo de quench}

Sus simulaciones usaron tres modelos de campo de fuerzas:
\begin{enumerate}
\item \textbf{mW} (Molinero-Moore): un átomo único por molécula, potencial 
Stillinger-Weber con función angular tetraédrica.
\item \textbf{SPC}: agua rígida con tres sitios de carga.
\item \textbf{F3C}: agua flexible con enlaces $\mathrm{O{-}H}$ ajustables.
\end{enumerate}

\noindent Para 192, 300 y 1000 moléculas de agua, se realizaron simulaciones de 
\emph{enfriamiento súbito} (\emph{quench}) desde dos temperaturas iniciales 
$\Tinit \in \{300, 370, 450\}$~\si{\kelvin} hasta $T_{\mathrm{bath}} = \SI{100}{\kelvin}$ 
(estado vidrioso).

\subsection{Resultados clave: el papel del hexámero}

Tras el quench, se midió la \emph{densidad de estados vibracionales} (DOS) en el 
rango de baja frecuencia $\nu \in [80, 160]\;\si{\per\centi\meter}$. Las 
observaciones fueron:

\begin{enumerate}
\item El agua templada desde $\Tinit \geq \SI{370}{\kelvin}$ presenta una DOS en 
este rango \emph{más cercana a la del hielo} que la templada desde $\Tinit \leq \SI{300}{\kelvin}$.
\item Este rango de frecuencias corresponde a las vibraciones colectivas del 
\emph{armazón} del \textbf{hexámero de agua} ---el complejo cíclico de seis 
moléculas--- considerado el núcleo cristalino primigenio para la formación de 
hielo $\mathrm{I_h}$.
\item La \emph{población} de hexámeros en agua a $\SI{370}{\kelvin}$ supera a la 
de agua a $\SI{300}{\kelvin}$.
\end{enumerate}

\noindent La conclusión: \emph{una mayor población de núcleos cristalinos 
pre-existentes en el agua caliente acelera drásticamente la transición de fase 
hacia el hielo durante el quench.} Esta es la interpretación moderna ``estructural'' 
del efecto Mpemba.

\subsection{Verificación por congelamiento de hexámero fijo}

Jin y Goddard realizaron además un experimento computacional decisivo: 
\emph{fijaron} las posiciones de las moléculas de un hexámero individual durante 
el quench desde $\SI{370}{\kelvin}$. El resultado: la correlación estructural con 
el hielo $\mathrm{I_h}$ después del quench fue \emph{mayor} que sin la 
restricción. Esto confirma que el hexámero actúa como un núcleo de cristalización.

\section{El paradigma polimérico de Naserifar y Goddard (2019)}

Naserifar y Goddard~\cite{naserifar2019} extendieron radicalmente el paradigma 
mediante simulaciones con un campo de fuerzas \emph{ab initio} de alta precisión 
(RexPoN), parametrizado contra CCSD(T) para los enlaces $\mathrm{O{-}H}$ y DFT 
para las interacciones de largo alcance.

\subsection{El descubrimiento: agua como polímero dinámico ramificado}

A diferencia del paradigma tradicional donde cada molécula de $\mathrm{H_2O}$ 
forma cuatro enlaces de hidrógeno fuertes (SHB, \emph{strong hydrogen bonds}, 
$r_{OO} \leq \SI{2.93}{\angstrom}$), las simulaciones RexPoN revelaron que:

\begin{itemize}
\item En el hielo: $\langle N_{\mathrm{SHB}} \rangle = 4.0$ (consistente con la 
estructura tetraédrica).
\item En el agua líquida a $\SI{273.5}{\kelvin}$ (recién fundida): 
$\langle N_{\mathrm{SHB}} \rangle$ \emph{cae abruptamente a 2.3}.
\item A $\SI{298}{\kelvin}$: $\langle N_{\mathrm{SHB}} \rangle = 2.1$.
\item A $\SI{400}{\kelvin}$: $\langle N_{\mathrm{SHB}} \rangle = 1.6$.
\end{itemize}

\noindent Esta caída implica que el agua líquida no es una ``red tetraédrica 
desordenada'' sino un \emph{conjunto de cadenas poliméricas ramificadas} en las 
que la mayoría de moléculas forma solo dos enlaces SHB.

\subsection{Estructura de la cadena polimérica}

A $\SI{298}{\kelvin}$, la cadena más larga contiene en promedio 
$\sim 151$ moléculas, organizadas en una topología ramificada con:

\begin{itemize}
\item Una \emph{columna vertebral} (\emph{backbone}) de \SIrange{30}{40}{} 
moléculas.
\item Múltiples ramas laterales de \SIrange{10}{20}{} moléculas.
\item Puntos terminales con un solo SHB.
\item Puntos de ramificación con tres SHB.
\end{itemize}

\noindent El tiempo de vida de un enlace SHB a $\SI{298}{\kelvin}$ es 
$\tau_{\mathrm{SHB}} = \SI{90.3}{\femto\second}$ ---aproximadamente 11 periodos 
de la vibración $\mathrm{O{-}H}$. Esto significa que la red polimérica se 
reconfigura constantemente, manteniendo coherencia estadística pero perdiendo 
identidad atómica en una escala de femtosegundos.

\subsection{Transición vítrea a 227 K}

El paradigma polimérico ofrece una interpretación natural del \emph{punto 
crítico del agua sobreenfriada} a $\SI{227}{\kelvin}$. Las simulaciones de 
Naserifar--Goddard predicen un salto del 15\% en $\tau_{\mathrm{SHB}}$ al cruzar 
esta temperatura. La interpretación: la transición no es de líquido-líquido sino 
de \emph{una transición vítrea de la red polimérica}.

\subsection{Conexión con el efecto Mpemba: el quench polimérico}

Naserifar y Goddard especulan~\cite{naserifar2019} que la diferencia entre el 
enfriamiento desde \SI{373}{\kelvin} y \SI{298}{\kelvin} se debe a las 
\emph{diferentes topologías de la red polimérica}:

\begin{itemize}
\item A $\SI{373}{\kelvin}$: cadenas cortas, $\langle N_{\mathrm{SHB}} \rangle = 1.7$, 
muchos núcleos potenciales para reorganización.
\item A $\SI{298}{\kelvin}$: cadenas largas (\SI{151}{} moléculas), atrapamiento 
en mínimos locales del paisaje energético.
\end{itemize}

\noindent Esta interpretación es consistente con la observación de 
Jin--Goddard~\cite{jin_goddard2015}: el agua caliente tiene mayor población de 
hexámeros porque sus cadenas más cortas se reconfiguran fácilmente en estructuras 
cíclicas, mientras que las cadenas largas del agua a temperatura ambiente quedan 
atrapadas en configuraciones extendidas que dificultan la cristalización.

\section{Síntesis del marco molecular}

Los tres marcos discutidos en este capítulo se complementan:

\begin{enumerate}
\item \textbf{Zhang \emph{et al.} 2014} aporta el modelo macroscópico de Fourier 
con memoria de hidrógeno y supersolidez de piel.
\item \textbf{Jin--Goddard 2015} identifica el hexámero de agua como núcleo 
microscópico y demuestra su mayor población en agua caliente.
\item \textbf{Naserifar--Goddard 2019} reemplaza la imagen tradicional de la red 
tetraédrica por una de polímeros ramificados dinámicos, y conecta la dinámica 
$\tau_{\mathrm{HB}}$ con el punto crítico del agua sobreenfriada.
\end{enumerate}

\noindent En el \cref{ch:implementacion} se describe cómo estos tres marcos están 
materializados en los módulos \texttt{water\_fourier} y \texttt{water\_md} de la 
suite \texttt{Mpemba-X v2}.

\begin{table}[h]
\centering
\caption{Parámetros del modelo molecular para el agua líquida según diferentes campos de fuerzas, comparados con valores experimentales.}
\label{tab:agua_params}
\small
\begin{tabular}{lcccc}
\toprule
\textbf{Propiedad} & \textbf{Exp.} & \textbf{TIP4P} & \textbf{mW} & \textbf{RexPoN} \\
\midrule
$T_{\mathrm{melt}}$ [\si{\kelvin}] & 273.15 & 232 & 274 & 273.3 \\
$\rho(\SI{298}{\kelvin})$ [\si{\gram\per\cubic\centi\meter}] & 0.997 & 0.999 & 0.994 & 0.997 \\
$\Delta H_{\mathrm{vap}}$ [\si{\kilo\calorie\per\mole}] & 10.52 & 10.65 & 11.05 & 10.36 \\
Entropía [\si{\joule\per\mole\per\kelvin}] & 69.9 & --- & --- & 68.4 \\
Constante dieléctrica & 78.4 & 53 & --- & 76.1 \\
$\langle N_{\mathrm{SHB}}\rangle$ ($r \leq \SI{2.93}{\angstrom}$) & 2.2\,$\pm$\,0.5 & 5.8 & 4.5 & 2.1 \\
\bottomrule
\end{tabular}
\end{table}
